\band{minthead}{Lecture 5 · Security Engineering}
\topic{Security landscape / CIA}
\k{Actor}: ผู้กระทำ; \k{Asset}: สิ่งที่ต้องปกป้อง เช่น บัญชี/รายงาน; \k{Attack surface}: จุดรับการโจมตี เช่น API/upload; \k{Trust boundary}: ขอบเขตความเชื่อถือ เช่น browser → server
\b{C: Confidentiality} ป้องกันการเห็นข้อมูลโดยไม่มีสิทธิ์; \b{I: Integrity} ป้องกันการแก้ไขโดยมิชอบ; \b{A: Availability} ผู้มีสิทธิ์ยังใช้ระบบได้
\k{Secure by design}: least privilege, defense in depth, fail safe/default deny, ออกแบบเรียบง่าย แบ่งส่วน ลดจุดโจมตี; ค่าเริ่มต้นปลอดภัยและยังใช้งานได้
\topic{Authentication ≠ Authorization}
\k{Authentication (AuthN)}: ยืนยันว่าเป็นใคร; \k{Authorization (AuthZ)}: ตรวจว่าทำได้ไหม ตามผู้ใช้ การกระทำ ทรัพยากรและ \r{state}; ตรวจที่ backend ทุก request
\k{RBAC}: สิทธิ์ตาม role; \k{ABAC}: ตามคุณลักษณะผู้ใช้/ทรัพยากร/บริบท; \k{ACL}: รายการสิทธิ์ต่อทรัพยากร; ตรวจ ownership เทียบตัวตนที่ server ยืนยัน \r{ไม่เชื่อคำอ้างว่าเป็นเจ้าของจาก client}
\k{Password}: slow hash + salt เฉพาะแต่ละรหัส; ไม่เก็บ/บันทึก plaintext. Base64 แค่แปลงรูปแบบ ไม่ใช่ encryption; encryption ถอดคืนด้วย key; hash ทางเดียว ใช้ตรวจเทียบ
\k{Basic}: ส่ง credentials แบบ Base64 ต้องใช้ HTTPS; \k{Digest}: ตอบด้วย hash/nonce เป็นวิธีเก่า ไม่แทน HTTPS
\k{Session}: server เก็บสถานะ ส่ง session ID ผ่าน cookie; หมดอายุ/logout แล้วใช้ต่อไม่ได้; current-user คืน profile ไม่คืน secret
\k{JWT}: token ที่ลงลายเซ็น \r{ไม่ได้เข้ารหัสเนื้อหาเสมอ}; ตรวจลายเซ็น/วันหมดอายุ; ต้องออกแบบเพิกถอนเอง; access อายุสั้น ส่วน refresh ต้องปกป้องและเพิกถอนได้
\k{API key}: secret ของแอป/บริการ; \k{OAuth2}: มอบสิทธิ์ให้แอปอื่น; \k{OIDC}: เพิ่มการยืนยันตัวตนบน OAuth2; \k{SSO}: sign-in ครั้งเดียวใช้หลายระบบ
\topic{OWASP / Vulnerability / Control}
{\smalltable\begin{tabularx}{\linewidth}{|p{10mm}|Y|}\hline
\rowcolor{lightrow}ช่องโหว่ & การป้องกัน / การทดสอบเชิงลบ\\\hline
Broken access & ตรวจ role/owner/state ที่ backend; เรียก API ข้ามคน/role → ปฏิเสธ ไม่รั่วข้อมูล\\\hline
Auth failure & error กลาง, rate limit, salted hash, expiry; ลอง wrong/inactive/logout แล้วเข้าตรง\\\hline
SQLi & parameterized query/ORM; allowlist ชื่อ field; \mbox{\code{' OR '1'='1}} ต้องเป็นข้อมูล ไม่ใช่ SQL\\\hline
XSS & encode output, sanitise rich text; ไม่ใช้ HTML ที่ไม่ปลอดภัย; payload ต้องไม่รัน\\\hline
CSRF & browser ส่ง cookie ข้าม site; ใช้ token/SameSite/Origin check; ปฏิเสธ forged write\\\hline
SSRF & server ดึง URL จาก input; allowlist ปลายทาง/บล็อก internal/metadata/loopback\\\hline
Upload & server ตรวจชนิด/ขนาด/เนื้อหา; เก็บแบบรันไม่ได้/ดาวน์โหลดปลอดภัย; script ปลอม PDF ต้องไม่รัน\\\hline
Config/secrets & HTTPS, secure cookie, CORS เคร่งครัด; ไม่เผย debug/.env; เปลี่ยน secret ที่รั่ว\\\hline
Deps & patch/audit packages; จำกัดสิทธิ์ DB/network; ปกป้อง backups\\\hline
\end{tabularx}}
\r{ซ่อนปุ่มไม่ใช่การควบคุมสิทธิ์}; frontend validation ช่วย UX; backend validation คุม trust boundary
\k{Status}: 400 ข้อมูลผิด; 401 ยังไม่ยืนยันตัวตน; 403 ไม่มีสิทธิ์; 404 ไม่พบ/ปิดบังตามกติกา; 409 ขัดแย้ง; 500 ข้อผิดพลาดกลาง; log correlation ID ไม่ใช่ password/token
\extra เมื่อปฏิเสธ ให้ตรวจ error/status และข้อมูลธุรกิจไม่เปลี่ยน; audit log ยังเกิดได้ตาม contract
\topic{คำอธิบายเสริม: ตัวอย่างตรวจสิทธิ์ผ่าน API ตรง}
{\smalltable\begin{tabularx}{\linewidth}{|p{13mm}|Y|}\hline
\rowcolor{lightrow}ทดสอบ & ผลที่ต้องตรวจตาม policy/contract\\\hline
AuthN/gate & valid ผ่าน; wrong/inactive/ไม่ login ถูกปฏิเสธ; password gate ต้องคุม API ตรง\\\hline
ผิด role & เรียก endpoint ตรงไม่ได้; ไม่มี business write\\\hline
ข้าม owner & A เปลี่ยน ID ไปขอข้อมูล B → ไม่เผยข้อมูล/จำนวน/ไฟล์\\\hline
logout/expiry & ใช้ session เดิมซ้ำไม่ได้ แม้เปิด URL/API ตรง\\\hline
input/duplicate & error ปลอดภัย; DB ไม่ค้างครึ่งชุด; ข้อมูลเดิมไม่เสีย\\\hline
log/config & ไม่พบ token/password/stack trace ใน response หรือ log ที่เปิดเผย\\\hline
\end{tabularx}}
\k{แนวตรวจ}: AuthN → AuthZ → validation → ลงมือ/บันทึก; ไม่เชื่อ role จาก client; ลำดับย่อยขึ้นกับข้อมูลที่ต้องใช้ตรวจสิทธิ์
\k{Defense in depth}: ต่อให้ไฟล์หลุด validation ต้องยังรันไม่ได้และอ่านข้าม owner ไม่ได้
\chapterband{minthead}{Lecture 6 · Project Initiation}
\b{Initiate → Plan → Execute → Close}; เป็นวงจรโครงการ ส่วน SDLC อยู่ภายในโครงการ
\k{Initiation 6}: goals, scope, deliverables, success criteria, stakeholders, resources → charter → อนุมัติ → วางแผน
\k{Goal}: ผลที่ต้องการ; \k{Deliverable}: สิ่งส่งมอบ; \k{Success criteria}: หลักฐานว่าสำเร็จ; launch = ส่งมอบ, landing = วัดผล
\k{SMART}: Specific ชัด, Measurable วัดได้, Attainable ทำได้, Relevant สอดคล้อง, Time-bound มีเส้นตาย; \g{ลดมัธยฐานเวลาตอบครั้งแรก 20\% ภายในภาคนี้ด้วย routing อัตโนมัติ}
\k{OKR}: Objective = ทิศทาง; KRs = ผลที่วัดได้; ท้าทายและเชื่อมองค์กร → แผนก → โครงการ; \r{KR ไม่ใช่รายการงาน}
\b{ROI=(Gain-Cost)/Cost ×100\%}; Gain = ผลตอบแทน, Cost = ต้นทุนทั้งหมดในช่วงเดียวกัน; พิจารณาคุณค่าที่ตีราคาไม่ได้แยก
\g{G=10,000; C=6,000+25×12=6,300 → ROI=58.7\%}
\k{Scope}: ตกลงสิ่งที่ทำ/ไม่ทำและข้อจำกัด; \k{Scope creep}: ขอบเขตเพิ่มโดยไม่ควบคุม ทั้งจากลูกค้าหรือทีม; ขอเปลี่ยน → วิเคราะห์ผล → เสนอทางเลือก → อนุมัติ → บันทึก
\k{Triangle}: scope/time/cost; เมื่อตรึงด้านหนึ่งต้องปรับด้านอื่น ไม่รับงานเพิ่มโดยไม่ประเมิน
\k{Stakeholders}: sponsor อนุมัติทุน, PM คุมแผน, ทีมทำ, ลูกค้าซื้อ, ผู้ใช้ใช้; เลือกวิธีดูแลตาม power/interest:
\begin{center}\begin{tikzpicture}[x=1mm,y=1mm,font=\fontsize{7.1}{7.4}\selectfont]
\fill[bluehead!28](5,0)rectangle(26,-6);\fill[orangehead!38](26,0)rectangle(47,-6);
\fill[lightrow](5,-6)rectangle(26,-12);\fill[minthead!35](26,-6)rectangle(47,-12);
\draw(5,0)rectangle(47,-12);\draw(26,0)--(26,-12);\draw(5,-6)--(47,-6);
\node at(15.5,-3){Keep satisfied};\node at(36.5,-3){Manage closely};
\node at(15.5,-9){Monitor};\node at(36.5,-9){Keep informed};
\draw[flow](3.2,-12)--(3.2,0);\node[rotate=90]at(1.3,-6){Power: Low → High};
\draw[flow](5,-13.5)--(47,-13.5);\node[below,inner sep=.5pt]at(26,-13.5){Interest: Low → High};
\end{tikzpicture}\end{center}
\k{Charter}: เหตุผล/เป้าหมาย/ขอบเขต/สิ่งส่งมอบ/เกณฑ์สำเร็จ/ทีม/ทรัพยากร/ทุน/กำหนดการ/สมมติฐาน/ความเสี่ยง/ผู้อนุมัติ

\extra เว็บที่ใช้ได้ = deliverable; ลดเวลารอ = goal; success criteria ระบุค่าตั้งต้น เป้าหมายและเส้นตายที่ตกลง
